% Appendix D - Future work: virtual laboratory with a chemical ontology
\section{Future Work: A Virtual Laboratory with a Chemical Ontology}\label{app:futurework}

This appendix describes a direction the team may pursue after the course. It is an idea,
not a part of the MVP or of the financial forecast, and nothing in it has been built.

\subsection{Why a Virtual Laboratory}
The quality of the patch depends on many choices: which photochromic molecule, which host material and
solvent, how thick the coating is, and whether other additives are present. Each combination has to be
tested for response time, fade time, and the effect of sunscreen (Appendix~\ref{app:photochromic}).
Testing every combination by hand would be slow and costly for a small team. A virtual laboratory
would record these experiments in a structured way, learn from them, and suggest which formulation to
test next.

Research groups call the physical counterpart a \emph{self-driving laboratory}, in which
automation and artificial intelligence plan and run experiments \cite{abolhasani2023,tom2024}.
Such laboratories are expensive, but a review of low-cost versions describes ``frugal twins,'' such as a
color-mixing laboratory that gives hands-on experience at low cost \cite{lo2024}. A color-based patch
test is close to that idea.

\subsection{Concept}
Figure~\ref{fig:vlab} shows three layers.

\begin{figure}[ht]
  \centering
  \begin{tikzpicture}[
      layer/.style={draw, rounded corners, minimum width=12.6cm, minimum height=1.35cm,
                    align=center, font=\small, text width=12.2cm},
      >=Stealth]
    \node[layer, fill=blue!6]  (a) at (0,0)    {\textbf{3. Virtual laboratory}\\Models predict color response and fade time from a formulation and conditions, and propose the next experiment};
    \node[layer, fill=blue!12] (b) at (0,-1.8) {\textbf{2. Experiment knowledge graph}\\Every patch test is stored with its formulation, conditions (UV index, temperature, SPF), measurements, and photographs};
    \node[layer, fill=blue!20] (c) at (0,-3.6) {\textbf{1. Chemical ontology}\\Photochromic molecules, polymers, solvents, UV filters, and their classes and roles};
    \draw[->, thick] (c) -- (b);
    \draw[->, thick] (b) -- (a);
  \end{tikzpicture}
  \caption{Proposed layers of a virtual laboratory for patch formulation (original concept).}
  \label{fig:vlab}
\end{figure}

\paragraph{Layer 1: chemical ontology.} An ontology gives every chemical a place in a classification
and a defined role, so that software can reason about them. ChEBI, an open ontology of chemical
entities, classifies entities by shared structure and by role, using relations such as ``is a'' and
``has role'' \cite{hastings2016}. We would reuse such an ontology where it covers our materials and add
our own classes for photochromic dyes, host materials, and UV filters where it does not. Whether
industrial dyes are covered would have to be checked.

\paragraph{Layer 2: experiment knowledge graph.} Following approaches already used in laboratories,
experiments are recorded as linked data: the ontology describes the materials and the steps, and each
run adds a record with its provenance \cite{bai2024}. Work on semantic models for wet-chemical synthesis
shows how free-text procedures, which contain vague phrases like ``room temperature,'' can be replaced by
structured steps that machines can read \cite{schilling2026}.

\paragraph{Layer 3: virtual laboratory.} With enough records, simple models can estimate how a new
formulation behaves, and the system can rank candidates before anyone mixes them. Real experiments
stay the final judge.

\subsection{Roadmap}
\begin{table}[ht]
\centering\small
\caption{Possible stages of the virtual laboratory}\label{tab:vlab-roadmap}
\begin{tabular}{@{}L{2.4cm}L{6.4cm}L{5.3cm}@{}}
  \toprule
  \textbf{Stage} & \textbf{Content} & \textbf{Needs}\\
  \midrule
  1. Structured log & Record every patch test in one table with fixed fields and photographs & Spreadsheet or small database; test protocol (Table~\ref{tab:photochromic-plan})\\
  \addlinespace
  2. Ontology and graph & Describe materials and tests with an ontology and store them as a graph & A chemist's advice; open tools\\
  \addlinespace
  3. Predictive models & Fit simple models of response and fade time to the data & Enough experiments; basic machine-learning skills\\
  \addlinespace
  4. Low-cost automation & A camera and UV lamp rig that photographs patches on a schedule & Small hardware budget; careful safety checks\\
  \bottomrule
\end{tabular}
\end{table}

\subsection{Limits and Risks}
\begin{itemize}[itemsep=0pt]
  \item The approach only pays off if the team runs many experiments; with a handful of tests, a
    spreadsheet is enough.
  \item Building and maintaining an ontology needs chemistry expertise that the team does not yet
    have in-house.
  \item Predictions from a model are only as good as the data, and a new material may behave
    differently from the ones already tested.
  \item The effort competes with work on the product and the market, so it should start only after
    the basic patch tests have given usable results.
\end{itemize}
